% Evaluation section, rewritten. Meant to be included from paper.tex with
% % Evaluation section, rewritten. Meant to be included from paper.tex with
% % Evaluation section, rewritten. Meant to be included from paper.tex with
% % Evaluation section, rewritten. Meant to be included from paper.tex with
% \input{evaluation_v2} at the point in the main text where the evaluation belongs.
% Requires three entries added to bibliography.bib: kirillov2019panoptic,
% luo2005coreference, krippendorff2004content.
%
% NOTE: the figures below are computed after normalising the gold file's
% "asymetrical" to the codebook's "asymmetrical" in the reciprocity column.
% Without that fix reciprocity reads 0.25 instead of 0.70 and every composite
% score drops. See the remark on label validation in "What the measure cannot see".
\section{Evaluating the annotation pipeline}
\label{sec:evaluation}

\subsection{What we are asking of the models}

Before running the codebook over some four hundred novels, we need to know what an
LLM annotation of a novel is actually worth. We therefore hand-annotated one novel,
Lucie Delarue-Mardrus's \emph{L'Ange et les pervers} (1930), following the same
codebook we give the models, and used it as a gold standard. The novel is a
deliberate choice rather than a convenient one: its protagonist, Marion / Mario de
Valdeclare, is presented as intersex and moves between sapphic salons and male
homosexual circles, so the novel exercises the parts of the grid we care most about
--- the \texttt{queer} value of \texttt{relation\_type}, collective and milieu
relationships, and desire that is coded rather than declared. The human annotation
records fourteen relationships, from the central bond between Marion and Laurette
Wells down to a two-line marriage recalled in the backstory.

One novel gives a first reading of how the pipeline behaves, not a corpus-level
estimate of its accuracy; everything below should be read as a pilot. Five models
were run on this novel through OpenRouter with an identical system prompt, codebook,
and JSON schema (Table~\ref{tab:models}).

\begin{table}[h]
  \centering
  \small
  \begin{tabular}{ll}
    \toprule
    Provider & Model \\
    \midrule
    DeepSeek  & V4 Pro \\
    Anthropic & Claude Sonnet 5 \\
    Anthropic & Claude Opus 5 \\
    OpenAI    & GPT-5.6 Luna Pro \\
    OpenAI    & GPT-5.6 Terra \\
    \bottomrule
  \end{tabular}
  \caption{The five models compared, all called through OpenRouter with an identical
  system prompt, codebook, and JSON schema.}
  \label{tab:models}
\end{table}

\subsection{Two questions, two measures}

Reading a novel with this grid means doing two rather different things. The
annotator first draws up a \emph{census}: which relationships in this novel are
worth an entry at all? Then, for each entry, they \emph{code} it, choosing one
label per closed field. A model can do one well and the other badly --- it can
find exactly the fourteen relationships a human found and misread the tone of
every one of them, or code impeccably three relationships nobody would have
listed. A single accuracy figure hides that difference, which is precisely the
difference we want to report. We therefore score the census and the coding
separately, then combine them.

The distinction is not ours. Content analysis has long separated
\emph{unitizing} --- deciding what the units of analysis are --- from
\emph{categorizing} them, and has insisted that the reliability of the first
cannot be folded into the reliability of the second
\cite{krippendorff2013content,krippendorff2016unitizing}. The measure that
addresses both at once is Mathet, Widlöcher and Métivier's $\gamma$, which
computes an alignment between annotators' units and scores agreement on that
alignment, with $\gamma_{cat}$ isolating the categorization component
\cite{mathet2015gamma,mathet2017gammacat}. We would use $\gamma$ if we could.
We cannot, because $\gamma$ assumes units positioned on a continuum --- spans of
text, stretches of recording --- and measures positional discrepancy as part of
its dissimilarity. Our units are not spans. A relationship is not located at a
place in the novel; it is identified by \emph{who is in it}, and its footprint
in the text is the whole book. What follows is therefore $\gamma$'s
architecture --- align first, then score categorization on the alignment ---
rebuilt on a set-identity criterion rather than a positional one.

\paragraph{Matching relationships.} Let $G$ be the set of relationships in the
human annotation of a novel and $P$ the set proposed by a model. Neither is an
ordered list, so nothing can be compared until the two sets are aligned. We
align them on the participants: each relationship $r$ carries an identifier
$\mathrm{id}(r)$ built from the characters involved ---
\texttt{laurette\_aimee} for the affair between Laurette Wells and Aimée de
Lagres, \texttt{pere\_mere} for Marion's parents --- recorded in both files in
the \texttt{id\_uniforme} column, with aliases resolved beforehand (Marion,
Mario, and Miss Marion Hervin are one character). Two relationships match when
their identifiers are equal:
%
\begin{align}
M = \{(g,p) \in G \times P \;:\; \mathrm{id}(g) = \mathrm{id}(p)\} .
\end{align}
%
In the vocabulary of coreference scoring this is CEAF under Luo's strictest
similarity, $\phi_1$, which credits an alignment only when the two sets are
identical \cite{luo2005coreference}. Luo sets $\phi_1$ aside as too severe for
coreference, where exact agreement on a mention chain is vanishingly rare; it is
usable here only because the participant lists are short and canonicalised
upstream. The gain is that the census score carries no similarity function and
no threshold --- it is a deterministic function of the two annotation files ---
and a disagreement about who is involved in a relationship is treated as what it
is: a disagreement about the novel, not a spelling problem.

\paragraph{Matching relationships.} Let $G$ be the set of relationships in the human
annotation of a novel and $P$ the set proposed by a model. Neither is an ordered
list, so nothing can be compared until the two sets are aligned. We align them on
the participants: each relationship $r$ carries an identifier $\mathrm{id}(r)$ built
from the characters involved --- \texttt{laurette\_aimee} for the affair between
Laurette Wells and Aimée de Lagres, \texttt{pere\_mere} for Marion's parents ---
recorded in both files in the \texttt{id\_uniforme} column, with aliases resolved
beforehand (Marion, Mario, and Miss Marion Hervin are one character). Two
relationships match when their identifiers are equal:
%
\begin{align}
M = \{(g,p) \in G \times P \;:\; \mathrm{id}(g) = \mathrm{id}(p)\} .
\end{align}
%
This is an exact, one-to-one join, with no similarity function and no threshold: the
census score is a deterministic function of the two annotation files, and a
disagreement about who is involved in a relationship is treated as what it is --- a
disagreement about the novel, not a spelling problem.

\paragraph{Recognition quality: did the model see the same relationships?} From $M$
we take the usual precision, recall, and their harmonic mean, computed over
relationships rather than tokens:
%
\begin{align}
\mathrm{P} = \frac{|M|}{|P|}, \qquad
\mathrm{R} = \frac{|M|}{|G|}, \qquad
\mathrm{RQ} = \frac{2\,\mathrm{P}\,\mathrm{R}}{\mathrm{P} + \mathrm{R}} .
\end{align}
%
Precision asks whether the relationships a model proposes are ones a human reader
would also have listed; recall asks how much of the novel's relational fabric it
leaves out. RQ says nothing at all about how well anything was described.

\paragraph{Coding quality: did it fill the grid the same way?} For a matched pair
$(g,p)$, we compare the closed fields one by one and take the proportion on which
the human and the model wrote the same label:
%
\begin{align}
a(g,p) &= \frac{1}{K} \sum_{k=1}^{K} \mathbf{1}\!\left[\, g_k = p_k \,\right] , \\
\mathrm{CQ} &= \frac{1}{|M|} \sum_{(g,p) \in M} a(g,p) .
\end{align}
%
Here $\mathbf{1}[\cdot]$ is 1 when the two labels are identical and 0 otherwise:
credit is strict, since the codebook offers no principled way to call
\texttt{decadent\_sublime} half-right when the annotator wrote
\texttt{moralized\_pathologized}. Every field weighs the same, $1/K$. The grid has
eleven closed fields, but the gold file predates \texttt{explicitness\_in\_plot},
added in schema v2, so this pilot scores $K = 10$: \texttt{relation\_type} and
\texttt{structure} (participants and shape); \texttt{status}, \texttt{reciprocity},
\texttt{exclusivity}, \texttt{adultery} (nature and rules of the bond);
\texttt{explicitness}, \texttt{centrality}, \texttt{framing}, \texttt{outcome}
(representation and narrative role). \texttt{confidence} is logged by the model but
is metadata about the annotation, not a property of the relationship, so it is not
scored. CQ is computed only on matched pairs: it measures description quality among
the relationships both readers agreed to list.

\paragraph{One number, and why it is a product.} Suppose we had scored the
whole annotation at once, weighting each matched relationship by how well it was
coded --- a precision and a recall in which a matched pair counts not as $1$ but
as $a(g,p)$. The harmonic mean of those two quantities is exactly
$\mathrm{RQ} \times \mathrm{CQ}$: the mean coding score factors out of the
numerator and cancels in the denominator. The product is therefore not a
weighting we chose but the single-figure score itself, read in its two parts.
%
\begin{align}
\mathrm{AQ} = \mathrm{RQ} \times \mathrm{CQ} .
\end{align}
%
This is the structure of panoptic quality in image segmentation, where a
recognition quality --- the same $F_1$ over matched objects --- multiplies a
segmentation quality that averages overlap on the matches
\cite{kirillov2019panoptic}. Ours differs in the similarity being averaged:
grid agreement in place of pixel overlap. The same two-stage shape governs the
scoring of template filling, from the MUC evaluations onward
\cite{chinchor1992muc4,du2021grit}, and Chen et al.\ \cite{chen2023unified} show
that this family of metrics --- match substructures, then normalise --- covers
most of structured-prediction evaluation. Naming that lineage is worth a
sentence because it locates a digital-humanities coding grid inside a
well-understood measurement problem rather than inventing a private one.

\subsection{What the pipeline gets right, and where it stops}

\begin{table}[h]
  \centering
  \small
  \begin{tabular}{lcccccc}
    \toprule
    Model & $|P|$ & P & R & RQ & CQ & AQ \\
    \midrule
    \textbf{Claude Opus 5} & 14 & 0.71 & 0.71 & 0.71 & 0.72 & \textbf{0.51} \\
    GPT-5.6 Terra          &  6 & 1.00 & 0.43 & 0.60 & 0.78 & 0.47 \\
    Claude Sonnet 5        &  6 & 1.00 & 0.43 & 0.60 & 0.73 & 0.44 \\
    GPT-5.6 Luna Pro       &  5 & 1.00 & 0.36 & 0.53 & 0.66 & 0.35 \\
    DeepSeek V4 Pro        &  7 & 0.86 & 0.43 & 0.57 & 0.47 & 0.27 \\
    \bottomrule
  \end{tabular}
  \caption{Census, coding, and combined score for the five models on
  \emph{L'Ange et les pervers}, against fourteen human-annotated relationships,
  sorted by AQ. $|P|$ is the number of relationships the model proposed.}
  \label{tab:model-results}
\end{table}

\paragraph{Two annotating temperaments.} The $|P|$ column of
Table~\ref{tab:model-results} is the most telling. Claude Opus 5 proposes exactly
fourteen relationships and gets ten of them right; the other four models propose five
to seven and are almost never wrong about the ones they propose. The second behaviour
is not caution about \emph{coding} but a narrower idea of what counts as a
relationship: those models list the named couples --- Laurette and Aimée, Aimée and
her ex-husband, Marion and Laurette --- and stop there. What they leave out is the
collective material: Laurette Wells surrounded by her favourites, the theatre milieu
around Ginette Lobre, the powdered young men of Julien Midalge's circle. For a study
of sapphic and homosexual sociability those are not decorative background; they are
the milieux the corpus was assembled to describe. A model with perfect precision that
returns only the foreground couples would quietly reshape the corpus into a series of
couples --- which is one reason precision alone is a poor guide to model choice here.

\paragraph{The grid has three layers, and they are not the codebook's three
families.} Averaging over fields hides the finding, so Table~\ref{tab:field-results}
reports agreement field by field.

\begin{table}[h]
  \centering
  \small
  \begin{tabular}{lcccccc}
    \toprule
    Field & Opus & Terra & Sonnet & Luna & DeepS. & mean \\
    \midrule
    \texttt{relation\_type} & 0.90 & 1.00 & 1.00 & 1.00 & 0.83 & 0.95 \\
    \texttt{structure}      & 0.90 & 1.00 & 0.67 & 1.00 & 1.00 & 0.91 \\
    \texttt{centrality}     & 0.90 & 0.83 & 1.00 & 0.80 & 0.67 & 0.84 \\
    \texttt{status}         & 0.70 & 1.00 & 0.83 & 0.80 & 0.83 & 0.83 \\
    \midrule
    \texttt{explicitness}   & 0.70 & 1.00 & 0.67 & 0.80 & 0.33 & 0.70 \\
    \texttt{reciprocity}    & 0.70 & 0.83 & 0.83 & 0.80 & 0.33 & 0.70 \\
    \texttt{outcome}        & 0.70 & 0.83 & 0.67 & 0.60 & 0.33 & 0.63 \\
    \midrule
    \texttt{exclusivity}    & 0.70 & 0.50 & 0.83 & 0.40 & 0.00 & 0.49 \\
    \texttt{adultery}       & 0.70 & 0.33 & 0.67 & 0.40 & 0.17 & 0.45 \\
    \texttt{framing}        & 0.30 & 0.50 & 0.17 & 0.00 & 0.17 & 0.23 \\
    \bottomrule
  \end{tabular}
  \caption{Agreement with the human annotation, field by field, on matched
  relationships only (10 pairs for Opus, 6 for Terra, Sonnet and DeepSeek, 5 for
  Luna). Fields are grouped by mean agreement, not by codebook family.}
  \label{tab:field-results}
\end{table}

The top group is the descriptive layer, and it is essentially solved. Once a model
and a human agree that a relationship exists, they agree on its gender configuration,
on the shape of the group, on the kind of bond, and on how much room the novel gives
it. The \texttt{queer} coding is the strongest case: every relationship involving
Marion / Mario is coded \texttt{queer} by all five models, without exception, which
is not a trivial result for a 1930 novel that never uses a clinical vocabulary for
its protagonist. In fact the only two \texttt{relation\_type} errors in the entire
pilot fall on the same relationship, Ginette Lobre and her theatre milieu, where the
annotator coded \texttt{mixed\_gender} and two models read a single-gender couple ---
that is, on the one relationship whose participants are a group rather than a person.

The middle group --- how openly the narration handles the relation, who is more
committed than whom, how the relation ends --- sits at roughly two thirds, usable
with care and worth reporting field by field rather than pooled.

The bottom group is where the pipeline fails, and the three fields have something in
common: they all ask for a judgment about norms. \texttt{exclusivity} and
\texttt{adultery} record the rules a relation operates under and whether they are
broken, and the codebook binds them by an explicit pairing rule (a relation that is
itself adulterous takes \texttt{adultery} together with \texttt{external\_adultery}).
The models do not apply that rule as a rule: pooling the matched pairs, they answer
\texttt{external\_adultery} four times where the annotator wrote \texttt{unclear} and
three times where she wrote \texttt{none}, and flatten \texttt{loyal} to \texttt{none}
four times. These are not fine judgments of the text but failures to treat the
codebook as a formal system, which suggests the remedy is a schema constraint rather
than a better prompt.

\texttt{framing}, the stance the novel takes towards a relationship, is the field
that resists most: 0.23 overall, against 0.11 for uniform guessing over its nine
labels. The error is systematic rather than random. \texttt{neutral\_ambivalent} is
the models' default answer, given fifteen times out of thirty-three, and
\texttt{redemptive\_spiritualized} --- which the annotator uses five times --- is
never once produced by any model. Where she reads spiritual rescue in Marion's
attachment to Laurette, three models read plain \texttt{positive} and the others
\texttt{decadent\_sublime} or \texttt{neutral\_ambivalent}; where she reads
\texttt{predatory\_or\_corrupting} in Laurette's hold over Aimée, the models spread
over \texttt{decadent\_sublime} and \texttt{neutral\_ambivalent}. The pattern is a
flattening: the models register that something intense is happening and decline to
say what the novel makes of it. The vocabulary of \emph{pervers} and of salvation
that gives the novel its title is exactly what the automatic pass does not recover.

That contrast \emph{is} the result we want to report: the grid contains a descriptive
layer that can be automated at corpus scale today and a normative layer that cannot,
and the second layer is the one that carries the interpretive claims of the project.

Two caveats keep these figures honest. Agreement is not chance-corrected, and the
marginals are skewed (\texttt{dyadic} dominates \texttt{structure}, \texttt{none}
dominates \texttt{adultery}), so the top of the table is flattered; a chance-corrected
coefficient such as Krippendorff's $\alpha$ \cite{krippendorff2004content} becomes
meaningful once a second human annotator gives us a ceiling to compare against. And
full-grid agreement is very rare: across thirty-three matched pairs, exactly two
relationships were coded correctly on all ten fields, both of them the same
relationship --- Aimée and her ex-husband, the most conventional bond in the book.

\paragraph{The easy relationships are not the interesting ones.} Reading the same
data by relationship rather than by model (Table~\ref{tab:relation-results}) sharpens
the point. The relationships every model finds and codes well are the peripheral,
legible ones: a marriage recalled in the backstory, a young man's one-sided
infatuation with Mario. The two relationships that organise the novel --- Marion and
Laurette, Laurette and Aimée --- are found by all five models and coded at 0.50,
half the grid. And the further one moves towards collective forms, the sharper the
drop: Laurette's circle of favourites is found once, Ginette's theatre milieu twice
and coded at 0.35, Midalge's homosexual circle never. The pipeline is most reliable
exactly where a scholar needs it least.

\begin{table}[h]
  \centering
  \small
  \begin{tabular}{llcc}
    \toprule
    Relationship & Type & Found & $\bar a$ \\
    \midrule
    Marion / Mario -- «~le petit brun~» & \texttt{queer}, episode & 5 & 0.84 \\
    Aimée -- comtesse Talliard       & \texttt{female\_female}    & 5 & 0.70 \\
    Marion / Mario -- Laurette Wells & \texttt{queer}, main       & 5 & 0.50 \\
    Laurette -- Aimée de Lagres      & \texttt{female\_female}    & 5 & 0.50 \\
    Aimée -- her ex-husband          & \texttt{female\_male}      & 4 & 0.95 \\
    Aimée -- the child's father      & \texttt{female\_male}      & 3 & 0.63 \\
    Marion's father -- Marion's mother & \texttt{female\_male}    & 2 & 0.90 \\
    Ginette Lobre -- theatre milieu  & \texttt{mixed\_gender}     & 2 & 0.35 \\
    Laurette -- Cecil Hampton        & \texttt{female\_male}      & 1 & 0.80 \\
    Laurette -- her favourites       & \texttt{hierarchy\_group}  & 1 & 0.70 \\
    \midrule
    Marion / Mario -- Aimée (seduction) & \texttt{queer}           & 0 & --- \\
    Marion / Mario -- Midalge's circle  & \texttt{queer}, milieu   & 0 & --- \\
    The uncle -- the aunt               & \texttt{female\_male}    & 0 & --- \\
    Marion -- the monk at Solesmes      & flagged \texttt{exclude} & 0 & --- \\
    \bottomrule
  \end{tabular}
  \caption{The fourteen human-annotated relationships, by number of models that found
  them (out of five) and mean coding agreement $\bar a$ among those models.}
  \label{tab:relation-results}
\end{table}

\paragraph{The models rarely decline.} The codebook gives every field an
\texttt{unclear} value, and the annotator uses it: some bonds in this novel genuinely
do not settle, particularly on \texttt{exclusivity} and \texttt{adultery} in the
milieu relationships. The models almost never do. Pooling matched pairs, where the
annotator wrote \texttt{unclear} the models committed to a substantive label in
twelve cases out of thirteen; in the other direction they abstained on eleven of the
three hundred and seventeen fields where she had committed. This asymmetry shapes
what a corpus-scale run would produce: not a cautious annotation with visible gaps,
but a complete-looking table whose confident cells include the ones a human reader
would have left open. Reporting it alongside the field agreements also stops
abstention from becoming a cheap way to look accurate.

\subsection{What the measure cannot see}

Matching on participants is what makes the census score exact, and it also fixes what
the score can register.

\emph{Collective relationships are the residual problem.} Alias resolution settles
naming for named characters; it does not decide whether Laurette and her favourites
are one \texttt{hierarchy\_group} relationship or six dyads. That is a scholarly
decision, not noise, and under an exact key a legitimate disagreement about it costs
a false positive and a false negative at once. The gap between the dyadic and
collective strata in Table~\ref{tab:relation-results} is therefore partly a measure
of how each model resolved a question the codebook leaves open. Fixing a convention
--- a canonical identifier per milieu, with members listed in a separate field --- is
the first change we plan for the next version of the grid.

\emph{Direction is invisible.} The grid stores participants as a list, not as roles,
while several fields are directional: \texttt{reciprocity} = \texttt{one\_sided} does
not say who desires whom, and \texttt{adultery} = \texttt{adultery\_one} does not say
which party breaks the rule. A model that inverts the seducer and the seduced scores
a perfect match. For a project whose object is the asymmetry of desire this is a
substantive gap, and the reason a \texttt{direction} field is the second planned
change.

\emph{A strict metric is only as good as the label vocabulary.} Exact matching makes
the score reproducible, but it also means that a variant spelling of a label counts
as a disagreement. One misspelling in the gold file --- \texttt{asymetrical} for the
codebook's \texttt{asymmetrical} --- silently drove agreement on
\texttt{reciprocity} from 0.70 down to 0.25, turning a working field into an apparent
failure. The scoring pipeline now validates every label in both files against the
codebook's declared vocabulary before any comparison is made, and we report this
because it is the kind of error that survives peer review unnoticed.

\emph{The gold standard is one novel, and the counts are small.} Ten matched pairs
for the best model means a single relationship moves CQ by ten points, and the
ordering of Opus, Terra, and Sonnet in Table~\ref{tab:model-results} should not be
read as a separation. One of the fourteen gold relationships (Marion and the monk at
Solesmes) carries the annotator's own \texttt{exclude} flag and is nonetheless
counted against the models here; treating it as an ignore region, as detection
benchmarks do, would raise every recall figure. And we have no human ceiling yet:
fields where two human annotators would agree at 0.6 cannot be used to convict a
model at 0.6, so double annotation of a subset is the measurement that would make
Table~\ref{tab:field-results} fully interpretable.

\subsection{What we carry forward}

On this pilot, Claude Opus 5 is the model we take to the rest of the corpus: it is
the only one that annotates the novel at the granularity the codebook asks for,
milieux included, and it pays for that coverage with a precision the others never
have to risk. The ranking rests on one novel and fourteen relationships and should be
revisited as the gold standard grows.

More important than the ranking is the shape of the result. The census and the
descriptive half of the grid --- who is involved, in what gender configuration, in
what shape of group, in what kind of bond, with how much narrative room --- are
reliable enough to run over four hundred novels and to support distributional claims.
The normative half --- what the relation owes to whom, and what the novel makes of it
--- is not, and we do not intend to publish corpus-wide figures on
\texttt{framing} on the strength of an automatic pass. The pipeline's role in this
project is accordingly a division of labour rather than a replacement: it produces
the census and the descriptive coding at scale, and it directs human reading towards
the relationships where the interpretive fields are actually at stake.
 at the point in the main text where the evaluation belongs.
% Requires three entries added to bibliography.bib: kirillov2019panoptic,
% luo2005coreference, krippendorff2004content.
%
% NOTE: the figures below are computed after normalising the gold file's
% "asymetrical" to the codebook's "asymmetrical" in the reciprocity column.
% Without that fix reciprocity reads 0.25 instead of 0.70 and every composite
% score drops. See the remark on label validation in "What the measure cannot see".
\section{Evaluating the annotation pipeline}
\label{sec:evaluation}

\subsection{What we are asking of the models}

Before running the codebook over some four hundred novels, we need to know what an
LLM annotation of a novel is actually worth. We therefore hand-annotated one novel,
Lucie Delarue-Mardrus's \emph{L'Ange et les pervers} (1930), following the same
codebook we give the models, and used it as a gold standard. The novel is a
deliberate choice rather than a convenient one: its protagonist, Marion / Mario de
Valdeclare, is presented as intersex and moves between sapphic salons and male
homosexual circles, so the novel exercises the parts of the grid we care most about
--- the \texttt{queer} value of \texttt{relation\_type}, collective and milieu
relationships, and desire that is coded rather than declared. The human annotation
records fourteen relationships, from the central bond between Marion and Laurette
Wells down to a two-line marriage recalled in the backstory.

One novel gives a first reading of how the pipeline behaves, not a corpus-level
estimate of its accuracy; everything below should be read as a pilot. Five models
were run on this novel through OpenRouter with an identical system prompt, codebook,
and JSON schema (Table~\ref{tab:models}).

\begin{table}[h]
  \centering
  \small
  \begin{tabular}{ll}
    \toprule
    Provider & Model \\
    \midrule
    DeepSeek  & V4 Pro \\
    Anthropic & Claude Sonnet 5 \\
    Anthropic & Claude Opus 5 \\
    OpenAI    & GPT-5.6 Luna Pro \\
    OpenAI    & GPT-5.6 Terra \\
    \bottomrule
  \end{tabular}
  \caption{The five models compared, all called through OpenRouter with an identical
  system prompt, codebook, and JSON schema.}
  \label{tab:models}
\end{table}

\subsection{Two questions, two measures}

Reading a novel with this grid means doing two rather different things. The
annotator first draws up a \emph{census}: which relationships in this novel are
worth an entry at all? Then, for each entry, they \emph{code} it, choosing one
label per closed field. A model can do one well and the other badly --- it can
find exactly the fourteen relationships a human found and misread the tone of
every one of them, or code impeccably three relationships nobody would have
listed. A single accuracy figure hides that difference, which is precisely the
difference we want to report. We therefore score the census and the coding
separately, then combine them.

The distinction is not ours. Content analysis has long separated
\emph{unitizing} --- deciding what the units of analysis are --- from
\emph{categorizing} them, and has insisted that the reliability of the first
cannot be folded into the reliability of the second
\cite{krippendorff2013content,krippendorff2016unitizing}. The measure that
addresses both at once is Mathet, Widlöcher and Métivier's $\gamma$, which
computes an alignment between annotators' units and scores agreement on that
alignment, with $\gamma_{cat}$ isolating the categorization component
\cite{mathet2015gamma,mathet2017gammacat}. We would use $\gamma$ if we could.
We cannot, because $\gamma$ assumes units positioned on a continuum --- spans of
text, stretches of recording --- and measures positional discrepancy as part of
its dissimilarity. Our units are not spans. A relationship is not located at a
place in the novel; it is identified by \emph{who is in it}, and its footprint
in the text is the whole book. What follows is therefore $\gamma$'s
architecture --- align first, then score categorization on the alignment ---
rebuilt on a set-identity criterion rather than a positional one.

\paragraph{Matching relationships.} Let $G$ be the set of relationships in the
human annotation of a novel and $P$ the set proposed by a model. Neither is an
ordered list, so nothing can be compared until the two sets are aligned. We
align them on the participants: each relationship $r$ carries an identifier
$\mathrm{id}(r)$ built from the characters involved ---
\texttt{laurette\_aimee} for the affair between Laurette Wells and Aimée de
Lagres, \texttt{pere\_mere} for Marion's parents --- recorded in both files in
the \texttt{id\_uniforme} column, with aliases resolved beforehand (Marion,
Mario, and Miss Marion Hervin are one character). Two relationships match when
their identifiers are equal:
%
\begin{align}
M = \{(g,p) \in G \times P \;:\; \mathrm{id}(g) = \mathrm{id}(p)\} .
\end{align}
%
In the vocabulary of coreference scoring this is CEAF under Luo's strictest
similarity, $\phi_1$, which credits an alignment only when the two sets are
identical \cite{luo2005coreference}. Luo sets $\phi_1$ aside as too severe for
coreference, where exact agreement on a mention chain is vanishingly rare; it is
usable here only because the participant lists are short and canonicalised
upstream. The gain is that the census score carries no similarity function and
no threshold --- it is a deterministic function of the two annotation files ---
and a disagreement about who is involved in a relationship is treated as what it
is: a disagreement about the novel, not a spelling problem.

\paragraph{Matching relationships.} Let $G$ be the set of relationships in the human
annotation of a novel and $P$ the set proposed by a model. Neither is an ordered
list, so nothing can be compared until the two sets are aligned. We align them on
the participants: each relationship $r$ carries an identifier $\mathrm{id}(r)$ built
from the characters involved --- \texttt{laurette\_aimee} for the affair between
Laurette Wells and Aimée de Lagres, \texttt{pere\_mere} for Marion's parents ---
recorded in both files in the \texttt{id\_uniforme} column, with aliases resolved
beforehand (Marion, Mario, and Miss Marion Hervin are one character). Two
relationships match when their identifiers are equal:
%
\begin{align}
M = \{(g,p) \in G \times P \;:\; \mathrm{id}(g) = \mathrm{id}(p)\} .
\end{align}
%
This is an exact, one-to-one join, with no similarity function and no threshold: the
census score is a deterministic function of the two annotation files, and a
disagreement about who is involved in a relationship is treated as what it is --- a
disagreement about the novel, not a spelling problem.

\paragraph{Recognition quality: did the model see the same relationships?} From $M$
we take the usual precision, recall, and their harmonic mean, computed over
relationships rather than tokens:
%
\begin{align}
\mathrm{P} = \frac{|M|}{|P|}, \qquad
\mathrm{R} = \frac{|M|}{|G|}, \qquad
\mathrm{RQ} = \frac{2\,\mathrm{P}\,\mathrm{R}}{\mathrm{P} + \mathrm{R}} .
\end{align}
%
Precision asks whether the relationships a model proposes are ones a human reader
would also have listed; recall asks how much of the novel's relational fabric it
leaves out. RQ says nothing at all about how well anything was described.

\paragraph{Coding quality: did it fill the grid the same way?} For a matched pair
$(g,p)$, we compare the closed fields one by one and take the proportion on which
the human and the model wrote the same label:
%
\begin{align}
a(g,p) &= \frac{1}{K} \sum_{k=1}^{K} \mathbf{1}\!\left[\, g_k = p_k \,\right] , \\
\mathrm{CQ} &= \frac{1}{|M|} \sum_{(g,p) \in M} a(g,p) .
\end{align}
%
Here $\mathbf{1}[\cdot]$ is 1 when the two labels are identical and 0 otherwise:
credit is strict, since the codebook offers no principled way to call
\texttt{decadent\_sublime} half-right when the annotator wrote
\texttt{moralized\_pathologized}. Every field weighs the same, $1/K$. The grid has
eleven closed fields, but the gold file predates \texttt{explicitness\_in\_plot},
added in schema v2, so this pilot scores $K = 10$: \texttt{relation\_type} and
\texttt{structure} (participants and shape); \texttt{status}, \texttt{reciprocity},
\texttt{exclusivity}, \texttt{adultery} (nature and rules of the bond);
\texttt{explicitness}, \texttt{centrality}, \texttt{framing}, \texttt{outcome}
(representation and narrative role). \texttt{confidence} is logged by the model but
is metadata about the annotation, not a property of the relationship, so it is not
scored. CQ is computed only on matched pairs: it measures description quality among
the relationships both readers agreed to list.

\paragraph{One number, and why it is a product.} Suppose we had scored the
whole annotation at once, weighting each matched relationship by how well it was
coded --- a precision and a recall in which a matched pair counts not as $1$ but
as $a(g,p)$. The harmonic mean of those two quantities is exactly
$\mathrm{RQ} \times \mathrm{CQ}$: the mean coding score factors out of the
numerator and cancels in the denominator. The product is therefore not a
weighting we chose but the single-figure score itself, read in its two parts.
%
\begin{align}
\mathrm{AQ} = \mathrm{RQ} \times \mathrm{CQ} .
\end{align}
%
This is the structure of panoptic quality in image segmentation, where a
recognition quality --- the same $F_1$ over matched objects --- multiplies a
segmentation quality that averages overlap on the matches
\cite{kirillov2019panoptic}. Ours differs in the similarity being averaged:
grid agreement in place of pixel overlap. The same two-stage shape governs the
scoring of template filling, from the MUC evaluations onward
\cite{chinchor1992muc4,du2021grit}, and Chen et al.\ \cite{chen2023unified} show
that this family of metrics --- match substructures, then normalise --- covers
most of structured-prediction evaluation. Naming that lineage is worth a
sentence because it locates a digital-humanities coding grid inside a
well-understood measurement problem rather than inventing a private one.

\subsection{What the pipeline gets right, and where it stops}

\begin{table}[h]
  \centering
  \small
  \begin{tabular}{lcccccc}
    \toprule
    Model & $|P|$ & P & R & RQ & CQ & AQ \\
    \midrule
    \textbf{Claude Opus 5} & 14 & 0.71 & 0.71 & 0.71 & 0.72 & \textbf{0.51} \\
    GPT-5.6 Terra          &  6 & 1.00 & 0.43 & 0.60 & 0.78 & 0.47 \\
    Claude Sonnet 5        &  6 & 1.00 & 0.43 & 0.60 & 0.73 & 0.44 \\
    GPT-5.6 Luna Pro       &  5 & 1.00 & 0.36 & 0.53 & 0.66 & 0.35 \\
    DeepSeek V4 Pro        &  7 & 0.86 & 0.43 & 0.57 & 0.47 & 0.27 \\
    \bottomrule
  \end{tabular}
  \caption{Census, coding, and combined score for the five models on
  \emph{L'Ange et les pervers}, against fourteen human-annotated relationships,
  sorted by AQ. $|P|$ is the number of relationships the model proposed.}
  \label{tab:model-results}
\end{table}

\paragraph{Two annotating temperaments.} The $|P|$ column of
Table~\ref{tab:model-results} is the most telling. Claude Opus 5 proposes exactly
fourteen relationships and gets ten of them right; the other four models propose five
to seven and are almost never wrong about the ones they propose. The second behaviour
is not caution about \emph{coding} but a narrower idea of what counts as a
relationship: those models list the named couples --- Laurette and Aimée, Aimée and
her ex-husband, Marion and Laurette --- and stop there. What they leave out is the
collective material: Laurette Wells surrounded by her favourites, the theatre milieu
around Ginette Lobre, the powdered young men of Julien Midalge's circle. For a study
of sapphic and homosexual sociability those are not decorative background; they are
the milieux the corpus was assembled to describe. A model with perfect precision that
returns only the foreground couples would quietly reshape the corpus into a series of
couples --- which is one reason precision alone is a poor guide to model choice here.

\paragraph{The grid has three layers, and they are not the codebook's three
families.} Averaging over fields hides the finding, so Table~\ref{tab:field-results}
reports agreement field by field.

\begin{table}[h]
  \centering
  \small
  \begin{tabular}{lcccccc}
    \toprule
    Field & Opus & Terra & Sonnet & Luna & DeepS. & mean \\
    \midrule
    \texttt{relation\_type} & 0.90 & 1.00 & 1.00 & 1.00 & 0.83 & 0.95 \\
    \texttt{structure}      & 0.90 & 1.00 & 0.67 & 1.00 & 1.00 & 0.91 \\
    \texttt{centrality}     & 0.90 & 0.83 & 1.00 & 0.80 & 0.67 & 0.84 \\
    \texttt{status}         & 0.70 & 1.00 & 0.83 & 0.80 & 0.83 & 0.83 \\
    \midrule
    \texttt{explicitness}   & 0.70 & 1.00 & 0.67 & 0.80 & 0.33 & 0.70 \\
    \texttt{reciprocity}    & 0.70 & 0.83 & 0.83 & 0.80 & 0.33 & 0.70 \\
    \texttt{outcome}        & 0.70 & 0.83 & 0.67 & 0.60 & 0.33 & 0.63 \\
    \midrule
    \texttt{exclusivity}    & 0.70 & 0.50 & 0.83 & 0.40 & 0.00 & 0.49 \\
    \texttt{adultery}       & 0.70 & 0.33 & 0.67 & 0.40 & 0.17 & 0.45 \\
    \texttt{framing}        & 0.30 & 0.50 & 0.17 & 0.00 & 0.17 & 0.23 \\
    \bottomrule
  \end{tabular}
  \caption{Agreement with the human annotation, field by field, on matched
  relationships only (10 pairs for Opus, 6 for Terra, Sonnet and DeepSeek, 5 for
  Luna). Fields are grouped by mean agreement, not by codebook family.}
  \label{tab:field-results}
\end{table}

The top group is the descriptive layer, and it is essentially solved. Once a model
and a human agree that a relationship exists, they agree on its gender configuration,
on the shape of the group, on the kind of bond, and on how much room the novel gives
it. The \texttt{queer} coding is the strongest case: every relationship involving
Marion / Mario is coded \texttt{queer} by all five models, without exception, which
is not a trivial result for a 1930 novel that never uses a clinical vocabulary for
its protagonist. In fact the only two \texttt{relation\_type} errors in the entire
pilot fall on the same relationship, Ginette Lobre and her theatre milieu, where the
annotator coded \texttt{mixed\_gender} and two models read a single-gender couple ---
that is, on the one relationship whose participants are a group rather than a person.

The middle group --- how openly the narration handles the relation, who is more
committed than whom, how the relation ends --- sits at roughly two thirds, usable
with care and worth reporting field by field rather than pooled.

The bottom group is where the pipeline fails, and the three fields have something in
common: they all ask for a judgment about norms. \texttt{exclusivity} and
\texttt{adultery} record the rules a relation operates under and whether they are
broken, and the codebook binds them by an explicit pairing rule (a relation that is
itself adulterous takes \texttt{adultery} together with \texttt{external\_adultery}).
The models do not apply that rule as a rule: pooling the matched pairs, they answer
\texttt{external\_adultery} four times where the annotator wrote \texttt{unclear} and
three times where she wrote \texttt{none}, and flatten \texttt{loyal} to \texttt{none}
four times. These are not fine judgments of the text but failures to treat the
codebook as a formal system, which suggests the remedy is a schema constraint rather
than a better prompt.

\texttt{framing}, the stance the novel takes towards a relationship, is the field
that resists most: 0.23 overall, against 0.11 for uniform guessing over its nine
labels. The error is systematic rather than random. \texttt{neutral\_ambivalent} is
the models' default answer, given fifteen times out of thirty-three, and
\texttt{redemptive\_spiritualized} --- which the annotator uses five times --- is
never once produced by any model. Where she reads spiritual rescue in Marion's
attachment to Laurette, three models read plain \texttt{positive} and the others
\texttt{decadent\_sublime} or \texttt{neutral\_ambivalent}; where she reads
\texttt{predatory\_or\_corrupting} in Laurette's hold over Aimée, the models spread
over \texttt{decadent\_sublime} and \texttt{neutral\_ambivalent}. The pattern is a
flattening: the models register that something intense is happening and decline to
say what the novel makes of it. The vocabulary of \emph{pervers} and of salvation
that gives the novel its title is exactly what the automatic pass does not recover.

That contrast \emph{is} the result we want to report: the grid contains a descriptive
layer that can be automated at corpus scale today and a normative layer that cannot,
and the second layer is the one that carries the interpretive claims of the project.

Two caveats keep these figures honest. Agreement is not chance-corrected, and the
marginals are skewed (\texttt{dyadic} dominates \texttt{structure}, \texttt{none}
dominates \texttt{adultery}), so the top of the table is flattered; a chance-corrected
coefficient such as Krippendorff's $\alpha$ \cite{krippendorff2004content} becomes
meaningful once a second human annotator gives us a ceiling to compare against. And
full-grid agreement is very rare: across thirty-three matched pairs, exactly two
relationships were coded correctly on all ten fields, both of them the same
relationship --- Aimée and her ex-husband, the most conventional bond in the book.

\paragraph{The easy relationships are not the interesting ones.} Reading the same
data by relationship rather than by model (Table~\ref{tab:relation-results}) sharpens
the point. The relationships every model finds and codes well are the peripheral,
legible ones: a marriage recalled in the backstory, a young man's one-sided
infatuation with Mario. The two relationships that organise the novel --- Marion and
Laurette, Laurette and Aimée --- are found by all five models and coded at 0.50,
half the grid. And the further one moves towards collective forms, the sharper the
drop: Laurette's circle of favourites is found once, Ginette's theatre milieu twice
and coded at 0.35, Midalge's homosexual circle never. The pipeline is most reliable
exactly where a scholar needs it least.

\begin{table}[h]
  \centering
  \small
  \begin{tabular}{llcc}
    \toprule
    Relationship & Type & Found & $\bar a$ \\
    \midrule
    Marion / Mario -- «~le petit brun~» & \texttt{queer}, episode & 5 & 0.84 \\
    Aimée -- comtesse Talliard       & \texttt{female\_female}    & 5 & 0.70 \\
    Marion / Mario -- Laurette Wells & \texttt{queer}, main       & 5 & 0.50 \\
    Laurette -- Aimée de Lagres      & \texttt{female\_female}    & 5 & 0.50 \\
    Aimée -- her ex-husband          & \texttt{female\_male}      & 4 & 0.95 \\
    Aimée -- the child's father      & \texttt{female\_male}      & 3 & 0.63 \\
    Marion's father -- Marion's mother & \texttt{female\_male}    & 2 & 0.90 \\
    Ginette Lobre -- theatre milieu  & \texttt{mixed\_gender}     & 2 & 0.35 \\
    Laurette -- Cecil Hampton        & \texttt{female\_male}      & 1 & 0.80 \\
    Laurette -- her favourites       & \texttt{hierarchy\_group}  & 1 & 0.70 \\
    \midrule
    Marion / Mario -- Aimée (seduction) & \texttt{queer}           & 0 & --- \\
    Marion / Mario -- Midalge's circle  & \texttt{queer}, milieu   & 0 & --- \\
    The uncle -- the aunt               & \texttt{female\_male}    & 0 & --- \\
    Marion -- the monk at Solesmes      & flagged \texttt{exclude} & 0 & --- \\
    \bottomrule
  \end{tabular}
  \caption{The fourteen human-annotated relationships, by number of models that found
  them (out of five) and mean coding agreement $\bar a$ among those models.}
  \label{tab:relation-results}
\end{table}

\paragraph{The models rarely decline.} The codebook gives every field an
\texttt{unclear} value, and the annotator uses it: some bonds in this novel genuinely
do not settle, particularly on \texttt{exclusivity} and \texttt{adultery} in the
milieu relationships. The models almost never do. Pooling matched pairs, where the
annotator wrote \texttt{unclear} the models committed to a substantive label in
twelve cases out of thirteen; in the other direction they abstained on eleven of the
three hundred and seventeen fields where she had committed. This asymmetry shapes
what a corpus-scale run would produce: not a cautious annotation with visible gaps,
but a complete-looking table whose confident cells include the ones a human reader
would have left open. Reporting it alongside the field agreements also stops
abstention from becoming a cheap way to look accurate.

\subsection{What the measure cannot see}

Matching on participants is what makes the census score exact, and it also fixes what
the score can register.

\emph{Collective relationships are the residual problem.} Alias resolution settles
naming for named characters; it does not decide whether Laurette and her favourites
are one \texttt{hierarchy\_group} relationship or six dyads. That is a scholarly
decision, not noise, and under an exact key a legitimate disagreement about it costs
a false positive and a false negative at once. The gap between the dyadic and
collective strata in Table~\ref{tab:relation-results} is therefore partly a measure
of how each model resolved a question the codebook leaves open. Fixing a convention
--- a canonical identifier per milieu, with members listed in a separate field --- is
the first change we plan for the next version of the grid.

\emph{Direction is invisible.} The grid stores participants as a list, not as roles,
while several fields are directional: \texttt{reciprocity} = \texttt{one\_sided} does
not say who desires whom, and \texttt{adultery} = \texttt{adultery\_one} does not say
which party breaks the rule. A model that inverts the seducer and the seduced scores
a perfect match. For a project whose object is the asymmetry of desire this is a
substantive gap, and the reason a \texttt{direction} field is the second planned
change.

\emph{A strict metric is only as good as the label vocabulary.} Exact matching makes
the score reproducible, but it also means that a variant spelling of a label counts
as a disagreement. One misspelling in the gold file --- \texttt{asymetrical} for the
codebook's \texttt{asymmetrical} --- silently drove agreement on
\texttt{reciprocity} from 0.70 down to 0.25, turning a working field into an apparent
failure. The scoring pipeline now validates every label in both files against the
codebook's declared vocabulary before any comparison is made, and we report this
because it is the kind of error that survives peer review unnoticed.

\emph{The gold standard is one novel, and the counts are small.} Ten matched pairs
for the best model means a single relationship moves CQ by ten points, and the
ordering of Opus, Terra, and Sonnet in Table~\ref{tab:model-results} should not be
read as a separation. One of the fourteen gold relationships (Marion and the monk at
Solesmes) carries the annotator's own \texttt{exclude} flag and is nonetheless
counted against the models here; treating it as an ignore region, as detection
benchmarks do, would raise every recall figure. And we have no human ceiling yet:
fields where two human annotators would agree at 0.6 cannot be used to convict a
model at 0.6, so double annotation of a subset is the measurement that would make
Table~\ref{tab:field-results} fully interpretable.

\subsection{What we carry forward}

On this pilot, Claude Opus 5 is the model we take to the rest of the corpus: it is
the only one that annotates the novel at the granularity the codebook asks for,
milieux included, and it pays for that coverage with a precision the others never
have to risk. The ranking rests on one novel and fourteen relationships and should be
revisited as the gold standard grows.

More important than the ranking is the shape of the result. The census and the
descriptive half of the grid --- who is involved, in what gender configuration, in
what shape of group, in what kind of bond, with how much narrative room --- are
reliable enough to run over four hundred novels and to support distributional claims.
The normative half --- what the relation owes to whom, and what the novel makes of it
--- is not, and we do not intend to publish corpus-wide figures on
\texttt{framing} on the strength of an automatic pass. The pipeline's role in this
project is accordingly a division of labour rather than a replacement: it produces
the census and the descriptive coding at scale, and it directs human reading towards
the relationships where the interpretive fields are actually at stake.
 at the point in the main text where the evaluation belongs.
% Requires three entries added to bibliography.bib: kirillov2019panoptic,
% luo2005coreference, krippendorff2004content.
%
% NOTE: the figures below are computed after normalising the gold file's
% "asymetrical" to the codebook's "asymmetrical" in the reciprocity column.
% Without that fix reciprocity reads 0.25 instead of 0.70 and every composite
% score drops. See the remark on label validation in "What the measure cannot see".
\section{Evaluating the annotation pipeline}
\label{sec:evaluation}

\subsection{What we are asking of the models}

Before running the codebook over some four hundred novels, we need to know what an
LLM annotation of a novel is actually worth. We therefore hand-annotated one novel,
Lucie Delarue-Mardrus's \emph{L'Ange et les pervers} (1930), following the same
codebook we give the models, and used it as a gold standard. The novel is a
deliberate choice rather than a convenient one: its protagonist, Marion / Mario de
Valdeclare, is presented as intersex and moves between sapphic salons and male
homosexual circles, so the novel exercises the parts of the grid we care most about
--- the \texttt{queer} value of \texttt{relation\_type}, collective and milieu
relationships, and desire that is coded rather than declared. The human annotation
records fourteen relationships, from the central bond between Marion and Laurette
Wells down to a two-line marriage recalled in the backstory.

One novel gives a first reading of how the pipeline behaves, not a corpus-level
estimate of its accuracy; everything below should be read as a pilot. Five models
were run on this novel through OpenRouter with an identical system prompt, codebook,
and JSON schema (Table~\ref{tab:models}).

\begin{table}[h]
  \centering
  \small
  \begin{tabular}{ll}
    \toprule
    Provider & Model \\
    \midrule
    DeepSeek  & V4 Pro \\
    Anthropic & Claude Sonnet 5 \\
    Anthropic & Claude Opus 5 \\
    OpenAI    & GPT-5.6 Luna Pro \\
    OpenAI    & GPT-5.6 Terra \\
    \bottomrule
  \end{tabular}
  \caption{The five models compared, all called through OpenRouter with an identical
  system prompt, codebook, and JSON schema.}
  \label{tab:models}
\end{table}

\subsection{Two questions, two measures}

Reading a novel with this grid means doing two rather different things. The
annotator first draws up a \emph{census}: which relationships in this novel are
worth an entry at all? Then, for each entry, they \emph{code} it, choosing one
label per closed field. A model can do one well and the other badly --- it can
find exactly the fourteen relationships a human found and misread the tone of
every one of them, or code impeccably three relationships nobody would have
listed. A single accuracy figure hides that difference, which is precisely the
difference we want to report. We therefore score the census and the coding
separately, then combine them.

The distinction is not ours. Content analysis has long separated
\emph{unitizing} --- deciding what the units of analysis are --- from
\emph{categorizing} them, and has insisted that the reliability of the first
cannot be folded into the reliability of the second
\cite{krippendorff2013content,krippendorff2016unitizing}. The measure that
addresses both at once is Mathet, Widlöcher and Métivier's $\gamma$, which
computes an alignment between annotators' units and scores agreement on that
alignment, with $\gamma_{cat}$ isolating the categorization component
\cite{mathet2015gamma,mathet2017gammacat}. We would use $\gamma$ if we could.
We cannot, because $\gamma$ assumes units positioned on a continuum --- spans of
text, stretches of recording --- and measures positional discrepancy as part of
its dissimilarity. Our units are not spans. A relationship is not located at a
place in the novel; it is identified by \emph{who is in it}, and its footprint
in the text is the whole book. What follows is therefore $\gamma$'s
architecture --- align first, then score categorization on the alignment ---
rebuilt on a set-identity criterion rather than a positional one.

\paragraph{Matching relationships.} Let $G$ be the set of relationships in the
human annotation of a novel and $P$ the set proposed by a model. Neither is an
ordered list, so nothing can be compared until the two sets are aligned. We
align them on the participants: each relationship $r$ carries an identifier
$\mathrm{id}(r)$ built from the characters involved ---
\texttt{laurette\_aimee} for the affair between Laurette Wells and Aimée de
Lagres, \texttt{pere\_mere} for Marion's parents --- recorded in both files in
the \texttt{id\_uniforme} column, with aliases resolved beforehand (Marion,
Mario, and Miss Marion Hervin are one character). Two relationships match when
their identifiers are equal:
%
\begin{align}
M = \{(g,p) \in G \times P \;:\; \mathrm{id}(g) = \mathrm{id}(p)\} .
\end{align}
%
In the vocabulary of coreference scoring this is CEAF under Luo's strictest
similarity, $\phi_1$, which credits an alignment only when the two sets are
identical \cite{luo2005coreference}. Luo sets $\phi_1$ aside as too severe for
coreference, where exact agreement on a mention chain is vanishingly rare; it is
usable here only because the participant lists are short and canonicalised
upstream. The gain is that the census score carries no similarity function and
no threshold --- it is a deterministic function of the two annotation files ---
and a disagreement about who is involved in a relationship is treated as what it
is: a disagreement about the novel, not a spelling problem.

\paragraph{Matching relationships.} Let $G$ be the set of relationships in the human
annotation of a novel and $P$ the set proposed by a model. Neither is an ordered
list, so nothing can be compared until the two sets are aligned. We align them on
the participants: each relationship $r$ carries an identifier $\mathrm{id}(r)$ built
from the characters involved --- \texttt{laurette\_aimee} for the affair between
Laurette Wells and Aimée de Lagres, \texttt{pere\_mere} for Marion's parents ---
recorded in both files in the \texttt{id\_uniforme} column, with aliases resolved
beforehand (Marion, Mario, and Miss Marion Hervin are one character). Two
relationships match when their identifiers are equal:
%
\begin{align}
M = \{(g,p) \in G \times P \;:\; \mathrm{id}(g) = \mathrm{id}(p)\} .
\end{align}
%
This is an exact, one-to-one join, with no similarity function and no threshold: the
census score is a deterministic function of the two annotation files, and a
disagreement about who is involved in a relationship is treated as what it is --- a
disagreement about the novel, not a spelling problem.

\paragraph{Recognition quality: did the model see the same relationships?} From $M$
we take the usual precision, recall, and their harmonic mean, computed over
relationships rather than tokens:
%
\begin{align}
\mathrm{P} = \frac{|M|}{|P|}, \qquad
\mathrm{R} = \frac{|M|}{|G|}, \qquad
\mathrm{RQ} = \frac{2\,\mathrm{P}\,\mathrm{R}}{\mathrm{P} + \mathrm{R}} .
\end{align}
%
Precision asks whether the relationships a model proposes are ones a human reader
would also have listed; recall asks how much of the novel's relational fabric it
leaves out. RQ says nothing at all about how well anything was described.

\paragraph{Coding quality: did it fill the grid the same way?} For a matched pair
$(g,p)$, we compare the closed fields one by one and take the proportion on which
the human and the model wrote the same label:
%
\begin{align}
a(g,p) &= \frac{1}{K} \sum_{k=1}^{K} \mathbf{1}\!\left[\, g_k = p_k \,\right] , \\
\mathrm{CQ} &= \frac{1}{|M|} \sum_{(g,p) \in M} a(g,p) .
\end{align}
%
Here $\mathbf{1}[\cdot]$ is 1 when the two labels are identical and 0 otherwise:
credit is strict, since the codebook offers no principled way to call
\texttt{decadent\_sublime} half-right when the annotator wrote
\texttt{moralized\_pathologized}. Every field weighs the same, $1/K$. The grid has
eleven closed fields, but the gold file predates \texttt{explicitness\_in\_plot},
added in schema v2, so this pilot scores $K = 10$: \texttt{relation\_type} and
\texttt{structure} (participants and shape); \texttt{status}, \texttt{reciprocity},
\texttt{exclusivity}, \texttt{adultery} (nature and rules of the bond);
\texttt{explicitness}, \texttt{centrality}, \texttt{framing}, \texttt{outcome}
(representation and narrative role). \texttt{confidence} is logged by the model but
is metadata about the annotation, not a property of the relationship, so it is not
scored. CQ is computed only on matched pairs: it measures description quality among
the relationships both readers agreed to list.

\paragraph{One number, and why it is a product.} Suppose we had scored the
whole annotation at once, weighting each matched relationship by how well it was
coded --- a precision and a recall in which a matched pair counts not as $1$ but
as $a(g,p)$. The harmonic mean of those two quantities is exactly
$\mathrm{RQ} \times \mathrm{CQ}$: the mean coding score factors out of the
numerator and cancels in the denominator. The product is therefore not a
weighting we chose but the single-figure score itself, read in its two parts.
%
\begin{align}
\mathrm{AQ} = \mathrm{RQ} \times \mathrm{CQ} .
\end{align}
%
This is the structure of panoptic quality in image segmentation, where a
recognition quality --- the same $F_1$ over matched objects --- multiplies a
segmentation quality that averages overlap on the matches
\cite{kirillov2019panoptic}. Ours differs in the similarity being averaged:
grid agreement in place of pixel overlap. The same two-stage shape governs the
scoring of template filling, from the MUC evaluations onward
\cite{chinchor1992muc4,du2021grit}, and Chen et al.\ \cite{chen2023unified} show
that this family of metrics --- match substructures, then normalise --- covers
most of structured-prediction evaluation. Naming that lineage is worth a
sentence because it locates a digital-humanities coding grid inside a
well-understood measurement problem rather than inventing a private one.

\subsection{What the pipeline gets right, and where it stops}

\begin{table}[h]
  \centering
  \small
  \begin{tabular}{lcccccc}
    \toprule
    Model & $|P|$ & P & R & RQ & CQ & AQ \\
    \midrule
    \textbf{Claude Opus 5} & 14 & 0.71 & 0.71 & 0.71 & 0.72 & \textbf{0.51} \\
    GPT-5.6 Terra          &  6 & 1.00 & 0.43 & 0.60 & 0.78 & 0.47 \\
    Claude Sonnet 5        &  6 & 1.00 & 0.43 & 0.60 & 0.73 & 0.44 \\
    GPT-5.6 Luna Pro       &  5 & 1.00 & 0.36 & 0.53 & 0.66 & 0.35 \\
    DeepSeek V4 Pro        &  7 & 0.86 & 0.43 & 0.57 & 0.47 & 0.27 \\
    \bottomrule
  \end{tabular}
  \caption{Census, coding, and combined score for the five models on
  \emph{L'Ange et les pervers}, against fourteen human-annotated relationships,
  sorted by AQ. $|P|$ is the number of relationships the model proposed.}
  \label{tab:model-results}
\end{table}

\paragraph{Two annotating temperaments.} The $|P|$ column of
Table~\ref{tab:model-results} is the most telling. Claude Opus 5 proposes exactly
fourteen relationships and gets ten of them right; the other four models propose five
to seven and are almost never wrong about the ones they propose. The second behaviour
is not caution about \emph{coding} but a narrower idea of what counts as a
relationship: those models list the named couples --- Laurette and Aimée, Aimée and
her ex-husband, Marion and Laurette --- and stop there. What they leave out is the
collective material: Laurette Wells surrounded by her favourites, the theatre milieu
around Ginette Lobre, the powdered young men of Julien Midalge's circle. For a study
of sapphic and homosexual sociability those are not decorative background; they are
the milieux the corpus was assembled to describe. A model with perfect precision that
returns only the foreground couples would quietly reshape the corpus into a series of
couples --- which is one reason precision alone is a poor guide to model choice here.

\paragraph{The grid has three layers, and they are not the codebook's three
families.} Averaging over fields hides the finding, so Table~\ref{tab:field-results}
reports agreement field by field.

\begin{table}[h]
  \centering
  \small
  \begin{tabular}{lcccccc}
    \toprule
    Field & Opus & Terra & Sonnet & Luna & DeepS. & mean \\
    \midrule
    \texttt{relation\_type} & 0.90 & 1.00 & 1.00 & 1.00 & 0.83 & 0.95 \\
    \texttt{structure}      & 0.90 & 1.00 & 0.67 & 1.00 & 1.00 & 0.91 \\
    \texttt{centrality}     & 0.90 & 0.83 & 1.00 & 0.80 & 0.67 & 0.84 \\
    \texttt{status}         & 0.70 & 1.00 & 0.83 & 0.80 & 0.83 & 0.83 \\
    \midrule
    \texttt{explicitness}   & 0.70 & 1.00 & 0.67 & 0.80 & 0.33 & 0.70 \\
    \texttt{reciprocity}    & 0.70 & 0.83 & 0.83 & 0.80 & 0.33 & 0.70 \\
    \texttt{outcome}        & 0.70 & 0.83 & 0.67 & 0.60 & 0.33 & 0.63 \\
    \midrule
    \texttt{exclusivity}    & 0.70 & 0.50 & 0.83 & 0.40 & 0.00 & 0.49 \\
    \texttt{adultery}       & 0.70 & 0.33 & 0.67 & 0.40 & 0.17 & 0.45 \\
    \texttt{framing}        & 0.30 & 0.50 & 0.17 & 0.00 & 0.17 & 0.23 \\
    \bottomrule
  \end{tabular}
  \caption{Agreement with the human annotation, field by field, on matched
  relationships only (10 pairs for Opus, 6 for Terra, Sonnet and DeepSeek, 5 for
  Luna). Fields are grouped by mean agreement, not by codebook family.}
  \label{tab:field-results}
\end{table}

The top group is the descriptive layer, and it is essentially solved. Once a model
and a human agree that a relationship exists, they agree on its gender configuration,
on the shape of the group, on the kind of bond, and on how much room the novel gives
it. The \texttt{queer} coding is the strongest case: every relationship involving
Marion / Mario is coded \texttt{queer} by all five models, without exception, which
is not a trivial result for a 1930 novel that never uses a clinical vocabulary for
its protagonist. In fact the only two \texttt{relation\_type} errors in the entire
pilot fall on the same relationship, Ginette Lobre and her theatre milieu, where the
annotator coded \texttt{mixed\_gender} and two models read a single-gender couple ---
that is, on the one relationship whose participants are a group rather than a person.

The middle group --- how openly the narration handles the relation, who is more
committed than whom, how the relation ends --- sits at roughly two thirds, usable
with care and worth reporting field by field rather than pooled.

The bottom group is where the pipeline fails, and the three fields have something in
common: they all ask for a judgment about norms. \texttt{exclusivity} and
\texttt{adultery} record the rules a relation operates under and whether they are
broken, and the codebook binds them by an explicit pairing rule (a relation that is
itself adulterous takes \texttt{adultery} together with \texttt{external\_adultery}).
The models do not apply that rule as a rule: pooling the matched pairs, they answer
\texttt{external\_adultery} four times where the annotator wrote \texttt{unclear} and
three times where she wrote \texttt{none}, and flatten \texttt{loyal} to \texttt{none}
four times. These are not fine judgments of the text but failures to treat the
codebook as a formal system, which suggests the remedy is a schema constraint rather
than a better prompt.

\texttt{framing}, the stance the novel takes towards a relationship, is the field
that resists most: 0.23 overall, against 0.11 for uniform guessing over its nine
labels. The error is systematic rather than random. \texttt{neutral\_ambivalent} is
the models' default answer, given fifteen times out of thirty-three, and
\texttt{redemptive\_spiritualized} --- which the annotator uses five times --- is
never once produced by any model. Where she reads spiritual rescue in Marion's
attachment to Laurette, three models read plain \texttt{positive} and the others
\texttt{decadent\_sublime} or \texttt{neutral\_ambivalent}; where she reads
\texttt{predatory\_or\_corrupting} in Laurette's hold over Aimée, the models spread
over \texttt{decadent\_sublime} and \texttt{neutral\_ambivalent}. The pattern is a
flattening: the models register that something intense is happening and decline to
say what the novel makes of it. The vocabulary of \emph{pervers} and of salvation
that gives the novel its title is exactly what the automatic pass does not recover.

That contrast \emph{is} the result we want to report: the grid contains a descriptive
layer that can be automated at corpus scale today and a normative layer that cannot,
and the second layer is the one that carries the interpretive claims of the project.

Two caveats keep these figures honest. Agreement is not chance-corrected, and the
marginals are skewed (\texttt{dyadic} dominates \texttt{structure}, \texttt{none}
dominates \texttt{adultery}), so the top of the table is flattered; a chance-corrected
coefficient such as Krippendorff's $\alpha$ \cite{krippendorff2004content} becomes
meaningful once a second human annotator gives us a ceiling to compare against. And
full-grid agreement is very rare: across thirty-three matched pairs, exactly two
relationships were coded correctly on all ten fields, both of them the same
relationship --- Aimée and her ex-husband, the most conventional bond in the book.

\paragraph{The easy relationships are not the interesting ones.} Reading the same
data by relationship rather than by model (Table~\ref{tab:relation-results}) sharpens
the point. The relationships every model finds and codes well are the peripheral,
legible ones: a marriage recalled in the backstory, a young man's one-sided
infatuation with Mario. The two relationships that organise the novel --- Marion and
Laurette, Laurette and Aimée --- are found by all five models and coded at 0.50,
half the grid. And the further one moves towards collective forms, the sharper the
drop: Laurette's circle of favourites is found once, Ginette's theatre milieu twice
and coded at 0.35, Midalge's homosexual circle never. The pipeline is most reliable
exactly where a scholar needs it least.

\begin{table}[h]
  \centering
  \small
  \begin{tabular}{llcc}
    \toprule
    Relationship & Type & Found & $\bar a$ \\
    \midrule
    Marion / Mario -- «~le petit brun~» & \texttt{queer}, episode & 5 & 0.84 \\
    Aimée -- comtesse Talliard       & \texttt{female\_female}    & 5 & 0.70 \\
    Marion / Mario -- Laurette Wells & \texttt{queer}, main       & 5 & 0.50 \\
    Laurette -- Aimée de Lagres      & \texttt{female\_female}    & 5 & 0.50 \\
    Aimée -- her ex-husband          & \texttt{female\_male}      & 4 & 0.95 \\
    Aimée -- the child's father      & \texttt{female\_male}      & 3 & 0.63 \\
    Marion's father -- Marion's mother & \texttt{female\_male}    & 2 & 0.90 \\
    Ginette Lobre -- theatre milieu  & \texttt{mixed\_gender}     & 2 & 0.35 \\
    Laurette -- Cecil Hampton        & \texttt{female\_male}      & 1 & 0.80 \\
    Laurette -- her favourites       & \texttt{hierarchy\_group}  & 1 & 0.70 \\
    \midrule
    Marion / Mario -- Aimée (seduction) & \texttt{queer}           & 0 & --- \\
    Marion / Mario -- Midalge's circle  & \texttt{queer}, milieu   & 0 & --- \\
    The uncle -- the aunt               & \texttt{female\_male}    & 0 & --- \\
    Marion -- the monk at Solesmes      & flagged \texttt{exclude} & 0 & --- \\
    \bottomrule
  \end{tabular}
  \caption{The fourteen human-annotated relationships, by number of models that found
  them (out of five) and mean coding agreement $\bar a$ among those models.}
  \label{tab:relation-results}
\end{table}

\paragraph{The models rarely decline.} The codebook gives every field an
\texttt{unclear} value, and the annotator uses it: some bonds in this novel genuinely
do not settle, particularly on \texttt{exclusivity} and \texttt{adultery} in the
milieu relationships. The models almost never do. Pooling matched pairs, where the
annotator wrote \texttt{unclear} the models committed to a substantive label in
twelve cases out of thirteen; in the other direction they abstained on eleven of the
three hundred and seventeen fields where she had committed. This asymmetry shapes
what a corpus-scale run would produce: not a cautious annotation with visible gaps,
but a complete-looking table whose confident cells include the ones a human reader
would have left open. Reporting it alongside the field agreements also stops
abstention from becoming a cheap way to look accurate.

\subsection{What the measure cannot see}

Matching on participants is what makes the census score exact, and it also fixes what
the score can register.

\emph{Collective relationships are the residual problem.} Alias resolution settles
naming for named characters; it does not decide whether Laurette and her favourites
are one \texttt{hierarchy\_group} relationship or six dyads. That is a scholarly
decision, not noise, and under an exact key a legitimate disagreement about it costs
a false positive and a false negative at once. The gap between the dyadic and
collective strata in Table~\ref{tab:relation-results} is therefore partly a measure
of how each model resolved a question the codebook leaves open. Fixing a convention
--- a canonical identifier per milieu, with members listed in a separate field --- is
the first change we plan for the next version of the grid.

\emph{Direction is invisible.} The grid stores participants as a list, not as roles,
while several fields are directional: \texttt{reciprocity} = \texttt{one\_sided} does
not say who desires whom, and \texttt{adultery} = \texttt{adultery\_one} does not say
which party breaks the rule. A model that inverts the seducer and the seduced scores
a perfect match. For a project whose object is the asymmetry of desire this is a
substantive gap, and the reason a \texttt{direction} field is the second planned
change.

\emph{A strict metric is only as good as the label vocabulary.} Exact matching makes
the score reproducible, but it also means that a variant spelling of a label counts
as a disagreement. One misspelling in the gold file --- \texttt{asymetrical} for the
codebook's \texttt{asymmetrical} --- silently drove agreement on
\texttt{reciprocity} from 0.70 down to 0.25, turning a working field into an apparent
failure. The scoring pipeline now validates every label in both files against the
codebook's declared vocabulary before any comparison is made, and we report this
because it is the kind of error that survives peer review unnoticed.

\emph{The gold standard is one novel, and the counts are small.} Ten matched pairs
for the best model means a single relationship moves CQ by ten points, and the
ordering of Opus, Terra, and Sonnet in Table~\ref{tab:model-results} should not be
read as a separation. One of the fourteen gold relationships (Marion and the monk at
Solesmes) carries the annotator's own \texttt{exclude} flag and is nonetheless
counted against the models here; treating it as an ignore region, as detection
benchmarks do, would raise every recall figure. And we have no human ceiling yet:
fields where two human annotators would agree at 0.6 cannot be used to convict a
model at 0.6, so double annotation of a subset is the measurement that would make
Table~\ref{tab:field-results} fully interpretable.

\subsection{What we carry forward}

On this pilot, Claude Opus 5 is the model we take to the rest of the corpus: it is
the only one that annotates the novel at the granularity the codebook asks for,
milieux included, and it pays for that coverage with a precision the others never
have to risk. The ranking rests on one novel and fourteen relationships and should be
revisited as the gold standard grows.

More important than the ranking is the shape of the result. The census and the
descriptive half of the grid --- who is involved, in what gender configuration, in
what shape of group, in what kind of bond, with how much narrative room --- are
reliable enough to run over four hundred novels and to support distributional claims.
The normative half --- what the relation owes to whom, and what the novel makes of it
--- is not, and we do not intend to publish corpus-wide figures on
\texttt{framing} on the strength of an automatic pass. The pipeline's role in this
project is accordingly a division of labour rather than a replacement: it produces
the census and the descriptive coding at scale, and it directs human reading towards
the relationships where the interpretive fields are actually at stake.
 at the point in the main text where the evaluation belongs.
% Requires three entries added to bibliography.bib: kirillov2019panoptic,
% luo2005coreference, krippendorff2004content.
%
% NOTE: the figures below are computed after normalising the gold file's
% "asymetrical" to the codebook's "asymmetrical" in the reciprocity column.
% Without that fix reciprocity reads 0.25 instead of 0.70 and every composite
% score drops. See the remark on label validation in "What the measure cannot see".
\section{Evaluating the annotation pipeline}
\label{sec:evaluation}

\subsection{What we are asking of the models}

Before running the codebook over some four hundred novels, we need to know what an
LLM annotation of a novel is actually worth. We therefore hand-annotated one novel,
Lucie Delarue-Mardrus's \emph{L'Ange et les pervers} (1930), following the same
codebook we give the models, and used it as a gold standard. The novel is a
deliberate choice rather than a convenient one: its protagonist, Marion / Mario de
Valdeclare, is presented as intersex and moves between sapphic salons and male
homosexual circles, so the novel exercises the parts of the grid we care most about
--- the \texttt{queer} value of \texttt{relation\_type}, collective and milieu
relationships, and desire that is coded rather than declared. The human annotation
records fourteen relationships, from the central bond between Marion and Laurette
Wells down to a two-line marriage recalled in the backstory.

One novel gives a first reading of how the pipeline behaves, not a corpus-level
estimate of its accuracy; everything below should be read as a pilot. Five models
were run on this novel through OpenRouter with an identical system prompt, codebook,
and JSON schema (Table~\ref{tab:models}).

\begin{table}[h]
  \centering
  \small
  \begin{tabular}{ll}
    \toprule
    Provider & Model \\
    \midrule
    DeepSeek  & V4 Pro \\
    Anthropic & Claude Sonnet 5 \\
    Anthropic & Claude Opus 5 \\
    OpenAI    & GPT-5.6 Luna Pro \\
    OpenAI    & GPT-5.6 Terra \\
    \bottomrule
  \end{tabular}
  \caption{The five models compared, all called through OpenRouter with an identical
  system prompt, codebook, and JSON schema.}
  \label{tab:models}
\end{table}

\subsection{Two questions, two measures}

Reading a novel with this grid means doing two rather different things. The
annotator first draws up a \emph{census}: which relationships in this novel are
worth an entry at all? Then, for each entry, they \emph{code} it, choosing one
label per closed field. A model can do one well and the other badly --- it can
find exactly the fourteen relationships a human found and misread the tone of
every one of them, or code impeccably three relationships nobody would have
listed. A single accuracy figure hides that difference, which is precisely the
difference we want to report. We therefore score the census and the coding
separately, then combine them.

The distinction is not ours. Content analysis has long separated
\emph{unitizing} --- deciding what the units of analysis are --- from
\emph{categorizing} them, and has insisted that the reliability of the first
cannot be folded into the reliability of the second
\cite{krippendorff2013content,krippendorff2016unitizing}. The measure that
addresses both at once is Mathet, Widlöcher and Métivier's $\gamma$, which
computes an alignment between annotators' units and scores agreement on that
alignment, with $\gamma_{cat}$ isolating the categorization component
\cite{mathet2015gamma,mathet2017gammacat}. We would use $\gamma$ if we could.
We cannot, because $\gamma$ assumes units positioned on a continuum --- spans of
text, stretches of recording --- and measures positional discrepancy as part of
its dissimilarity. Our units are not spans. A relationship is not located at a
place in the novel; it is identified by \emph{who is in it}, and its footprint
in the text is the whole book. What follows is therefore $\gamma$'s
architecture --- align first, then score categorization on the alignment ---
rebuilt on a set-identity criterion rather than a positional one.

\paragraph{Matching relationships.} Let $G$ be the set of relationships in the
human annotation of a novel and $P$ the set proposed by a model. Neither is an
ordered list, so nothing can be compared until the two sets are aligned. We
align them on the participants: each relationship $r$ carries an identifier
$\mathrm{id}(r)$ built from the characters involved ---
\texttt{laurette\_aimee} for the affair between Laurette Wells and Aimée de
Lagres, \texttt{pere\_mere} for Marion's parents --- recorded in both files in
the \texttt{id\_uniforme} column, with aliases resolved beforehand (Marion,
Mario, and Miss Marion Hervin are one character). Two relationships match when
their identifiers are equal:
%
\begin{align}
M = \{(g,p) \in G \times P \;:\; \mathrm{id}(g) = \mathrm{id}(p)\} .
\end{align}
%
In the vocabulary of coreference scoring this is CEAF under Luo's strictest
similarity, $\phi_1$, which credits an alignment only when the two sets are
identical \cite{luo2005coreference}. Luo sets $\phi_1$ aside as too severe for
coreference, where exact agreement on a mention chain is vanishingly rare; it is
usable here only because the participant lists are short and canonicalised
upstream. The gain is that the census score carries no similarity function and
no threshold --- it is a deterministic function of the two annotation files ---
and a disagreement about who is involved in a relationship is treated as what it
is: a disagreement about the novel, not a spelling problem.

\paragraph{Matching relationships.} Let $G$ be the set of relationships in the human
annotation of a novel and $P$ the set proposed by a model. Neither is an ordered
list, so nothing can be compared until the two sets are aligned. We align them on
the participants: each relationship $r$ carries an identifier $\mathrm{id}(r)$ built
from the characters involved --- \texttt{laurette\_aimee} for the affair between
Laurette Wells and Aimée de Lagres, \texttt{pere\_mere} for Marion's parents ---
recorded in both files in the \texttt{id\_uniforme} column, with aliases resolved
beforehand (Marion, Mario, and Miss Marion Hervin are one character). Two
relationships match when their identifiers are equal:
%
\begin{align}
M = \{(g,p) \in G \times P \;:\; \mathrm{id}(g) = \mathrm{id}(p)\} .
\end{align}
%
This is an exact, one-to-one join, with no similarity function and no threshold: the
census score is a deterministic function of the two annotation files, and a
disagreement about who is involved in a relationship is treated as what it is --- a
disagreement about the novel, not a spelling problem.

\paragraph{Recognition quality: did the model see the same relationships?} From $M$
we take the usual precision, recall, and their harmonic mean, computed over
relationships rather than tokens:
%
\begin{align}
\mathrm{P} = \frac{|M|}{|P|}, \qquad
\mathrm{R} = \frac{|M|}{|G|}, \qquad
\mathrm{RQ} = \frac{2\,\mathrm{P}\,\mathrm{R}}{\mathrm{P} + \mathrm{R}} .
\end{align}
%
Precision asks whether the relationships a model proposes are ones a human reader
would also have listed; recall asks how much of the novel's relational fabric it
leaves out. RQ says nothing at all about how well anything was described.

\paragraph{Coding quality: did it fill the grid the same way?} For a matched pair
$(g,p)$, we compare the closed fields one by one and take the proportion on which
the human and the model wrote the same label:
%
\begin{align}
a(g,p) &= \frac{1}{K} \sum_{k=1}^{K} \mathbf{1}\!\left[\, g_k = p_k \,\right] , \\
\mathrm{CQ} &= \frac{1}{|M|} \sum_{(g,p) \in M} a(g,p) .
\end{align}
%
Here $\mathbf{1}[\cdot]$ is 1 when the two labels are identical and 0 otherwise:
credit is strict, since the codebook offers no principled way to call
\texttt{decadent\_sublime} half-right when the annotator wrote
\texttt{moralized\_pathologized}. Every field weighs the same, $1/K$. The grid has
eleven closed fields, but the gold file predates \texttt{explicitness\_in\_plot},
added in schema v2, so this pilot scores $K = 10$: \texttt{relation\_type} and
\texttt{structure} (participants and shape); \texttt{status}, \texttt{reciprocity},
\texttt{exclusivity}, \texttt{adultery} (nature and rules of the bond);
\texttt{explicitness}, \texttt{centrality}, \texttt{framing}, \texttt{outcome}
(representation and narrative role). \texttt{confidence} is logged by the model but
is metadata about the annotation, not a property of the relationship, so it is not
scored. CQ is computed only on matched pairs: it measures description quality among
the relationships both readers agreed to list.

\paragraph{One number, and why it is a product.} Suppose we had scored the
whole annotation at once, weighting each matched relationship by how well it was
coded --- a precision and a recall in which a matched pair counts not as $1$ but
as $a(g,p)$. The harmonic mean of those two quantities is exactly
$\mathrm{RQ} \times \mathrm{CQ}$: the mean coding score factors out of the
numerator and cancels in the denominator. The product is therefore not a
weighting we chose but the single-figure score itself, read in its two parts.
%
\begin{align}
\mathrm{AQ} = \mathrm{RQ} \times \mathrm{CQ} .
\end{align}
%
This is the structure of panoptic quality in image segmentation, where a
recognition quality --- the same $F_1$ over matched objects --- multiplies a
segmentation quality that averages overlap on the matches
\cite{kirillov2019panoptic}. Ours differs in the similarity being averaged:
grid agreement in place of pixel overlap. The same two-stage shape governs the
scoring of template filling, from the MUC evaluations onward
\cite{chinchor1992muc4,du2021grit}, and Chen et al.\ \cite{chen2023unified} show
that this family of metrics --- match substructures, then normalise --- covers
most of structured-prediction evaluation. Naming that lineage is worth a
sentence because it locates a digital-humanities coding grid inside a
well-understood measurement problem rather than inventing a private one.

\subsection{What the pipeline gets right, and where it stops}

\begin{table}[h]
  \centering
  \small
  \begin{tabular}{lcccccc}
    \toprule
    Model & $|P|$ & P & R & RQ & CQ & AQ \\
    \midrule
    \textbf{Claude Opus 5} & 14 & 0.71 & 0.71 & 0.71 & 0.72 & \textbf{0.51} \\
    GPT-5.6 Terra          &  6 & 1.00 & 0.43 & 0.60 & 0.78 & 0.47 \\
    Claude Sonnet 5        &  6 & 1.00 & 0.43 & 0.60 & 0.73 & 0.44 \\
    GPT-5.6 Luna Pro       &  5 & 1.00 & 0.36 & 0.53 & 0.66 & 0.35 \\
    DeepSeek V4 Pro        &  7 & 0.86 & 0.43 & 0.57 & 0.47 & 0.27 \\
    \bottomrule
  \end{tabular}
  \caption{Census, coding, and combined score for the five models on
  \emph{L'Ange et les pervers}, against fourteen human-annotated relationships,
  sorted by AQ. $|P|$ is the number of relationships the model proposed.}
  \label{tab:model-results}
\end{table}

\paragraph{Two annotating temperaments.} The $|P|$ column of
Table~\ref{tab:model-results} is the most telling. Claude Opus 5 proposes exactly
fourteen relationships and gets ten of them right; the other four models propose five
to seven and are almost never wrong about the ones they propose. The second behaviour
is not caution about \emph{coding} but a narrower idea of what counts as a
relationship: those models list the named couples --- Laurette and Aimée, Aimée and
her ex-husband, Marion and Laurette --- and stop there. What they leave out is the
collective material: Laurette Wells surrounded by her favourites, the theatre milieu
around Ginette Lobre, the powdered young men of Julien Midalge's circle. For a study
of sapphic and homosexual sociability those are not decorative background; they are
the milieux the corpus was assembled to describe. A model with perfect precision that
returns only the foreground couples would quietly reshape the corpus into a series of
couples --- which is one reason precision alone is a poor guide to model choice here.

\paragraph{The grid has three layers, and they are not the codebook's three
families.} Averaging over fields hides the finding, so Table~\ref{tab:field-results}
reports agreement field by field.

\begin{table}[h]
  \centering
  \small
  \begin{tabular}{lcccccc}
    \toprule
    Field & Opus & Terra & Sonnet & Luna & DeepS. & mean \\
    \midrule
    \texttt{relation\_type} & 0.90 & 1.00 & 1.00 & 1.00 & 0.83 & 0.95 \\
    \texttt{structure}      & 0.90 & 1.00 & 0.67 & 1.00 & 1.00 & 0.91 \\
    \texttt{centrality}     & 0.90 & 0.83 & 1.00 & 0.80 & 0.67 & 0.84 \\
    \texttt{status}         & 0.70 & 1.00 & 0.83 & 0.80 & 0.83 & 0.83 \\
    \midrule
    \texttt{explicitness}   & 0.70 & 1.00 & 0.67 & 0.80 & 0.33 & 0.70 \\
    \texttt{reciprocity}    & 0.70 & 0.83 & 0.83 & 0.80 & 0.33 & 0.70 \\
    \texttt{outcome}        & 0.70 & 0.83 & 0.67 & 0.60 & 0.33 & 0.63 \\
    \midrule
    \texttt{exclusivity}    & 0.70 & 0.50 & 0.83 & 0.40 & 0.00 & 0.49 \\
    \texttt{adultery}       & 0.70 & 0.33 & 0.67 & 0.40 & 0.17 & 0.45 \\
    \texttt{framing}        & 0.30 & 0.50 & 0.17 & 0.00 & 0.17 & 0.23 \\
    \bottomrule
  \end{tabular}
  \caption{Agreement with the human annotation, field by field, on matched
  relationships only (10 pairs for Opus, 6 for Terra, Sonnet and DeepSeek, 5 for
  Luna). Fields are grouped by mean agreement, not by codebook family.}
  \label{tab:field-results}
\end{table}

The top group is the descriptive layer, and it is essentially solved. Once a model
and a human agree that a relationship exists, they agree on its gender configuration,
on the shape of the group, on the kind of bond, and on how much room the novel gives
it. The \texttt{queer} coding is the strongest case: every relationship involving
Marion / Mario is coded \texttt{queer} by all five models, without exception, which
is not a trivial result for a 1930 novel that never uses a clinical vocabulary for
its protagonist. In fact the only two \texttt{relation\_type} errors in the entire
pilot fall on the same relationship, Ginette Lobre and her theatre milieu, where the
annotator coded \texttt{mixed\_gender} and two models read a single-gender couple ---
that is, on the one relationship whose participants are a group rather than a person.

The middle group --- how openly the narration handles the relation, who is more
committed than whom, how the relation ends --- sits at roughly two thirds, usable
with care and worth reporting field by field rather than pooled.

The bottom group is where the pipeline fails, and the three fields have something in
common: they all ask for a judgment about norms. \texttt{exclusivity} and
\texttt{adultery} record the rules a relation operates under and whether they are
broken, and the codebook binds them by an explicit pairing rule (a relation that is
itself adulterous takes \texttt{adultery} together with \texttt{external\_adultery}).
The models do not apply that rule as a rule: pooling the matched pairs, they answer
\texttt{external\_adultery} four times where the annotator wrote \texttt{unclear} and
three times where she wrote \texttt{none}, and flatten \texttt{loyal} to \texttt{none}
four times. These are not fine judgments of the text but failures to treat the
codebook as a formal system, which suggests the remedy is a schema constraint rather
than a better prompt.

\texttt{framing}, the stance the novel takes towards a relationship, is the field
that resists most: 0.23 overall, against 0.11 for uniform guessing over its nine
labels. The error is systematic rather than random. \texttt{neutral\_ambivalent} is
the models' default answer, given fifteen times out of thirty-three, and
\texttt{redemptive\_spiritualized} --- which the annotator uses five times --- is
never once produced by any model. Where she reads spiritual rescue in Marion's
attachment to Laurette, three models read plain \texttt{positive} and the others
\texttt{decadent\_sublime} or \texttt{neutral\_ambivalent}; where she reads
\texttt{predatory\_or\_corrupting} in Laurette's hold over Aimée, the models spread
over \texttt{decadent\_sublime} and \texttt{neutral\_ambivalent}. The pattern is a
flattening: the models register that something intense is happening and decline to
say what the novel makes of it. The vocabulary of \emph{pervers} and of salvation
that gives the novel its title is exactly what the automatic pass does not recover.

That contrast \emph{is} the result we want to report: the grid contains a descriptive
layer that can be automated at corpus scale today and a normative layer that cannot,
and the second layer is the one that carries the interpretive claims of the project.

Two caveats keep these figures honest. Agreement is not chance-corrected, and the
marginals are skewed (\texttt{dyadic} dominates \texttt{structure}, \texttt{none}
dominates \texttt{adultery}), so the top of the table is flattered; a chance-corrected
coefficient such as Krippendorff's $\alpha$ \cite{krippendorff2004content} becomes
meaningful once a second human annotator gives us a ceiling to compare against. And
full-grid agreement is very rare: across thirty-three matched pairs, exactly two
relationships were coded correctly on all ten fields, both of them the same
relationship --- Aimée and her ex-husband, the most conventional bond in the book.

\paragraph{The easy relationships are not the interesting ones.} Reading the same
data by relationship rather than by model (Table~\ref{tab:relation-results}) sharpens
the point. The relationships every model finds and codes well are the peripheral,
legible ones: a marriage recalled in the backstory, a young man's one-sided
infatuation with Mario. The two relationships that organise the novel --- Marion and
Laurette, Laurette and Aimée --- are found by all five models and coded at 0.50,
half the grid. And the further one moves towards collective forms, the sharper the
drop: Laurette's circle of favourites is found once, Ginette's theatre milieu twice
and coded at 0.35, Midalge's homosexual circle never. The pipeline is most reliable
exactly where a scholar needs it least.

\begin{table}[h]
  \centering
  \small
  \begin{tabular}{llcc}
    \toprule
    Relationship & Type & Found & $\bar a$ \\
    \midrule
    Marion / Mario -- «~le petit brun~» & \texttt{queer}, episode & 5 & 0.84 \\
    Aimée -- comtesse Talliard       & \texttt{female\_female}    & 5 & 0.70 \\
    Marion / Mario -- Laurette Wells & \texttt{queer}, main       & 5 & 0.50 \\
    Laurette -- Aimée de Lagres      & \texttt{female\_female}    & 5 & 0.50 \\
    Aimée -- her ex-husband          & \texttt{female\_male}      & 4 & 0.95 \\
    Aimée -- the child's father      & \texttt{female\_male}      & 3 & 0.63 \\
    Marion's father -- Marion's mother & \texttt{female\_male}    & 2 & 0.90 \\
    Ginette Lobre -- theatre milieu  & \texttt{mixed\_gender}     & 2 & 0.35 \\
    Laurette -- Cecil Hampton        & \texttt{female\_male}      & 1 & 0.80 \\
    Laurette -- her favourites       & \texttt{hierarchy\_group}  & 1 & 0.70 \\
    \midrule
    Marion / Mario -- Aimée (seduction) & \texttt{queer}           & 0 & --- \\
    Marion / Mario -- Midalge's circle  & \texttt{queer}, milieu   & 0 & --- \\
    The uncle -- the aunt               & \texttt{female\_male}    & 0 & --- \\
    Marion -- the monk at Solesmes      & flagged \texttt{exclude} & 0 & --- \\
    \bottomrule
  \end{tabular}
  \caption{The fourteen human-annotated relationships, by number of models that found
  them (out of five) and mean coding agreement $\bar a$ among those models.}
  \label{tab:relation-results}
\end{table}

\paragraph{The models rarely decline.} The codebook gives every field an
\texttt{unclear} value, and the annotator uses it: some bonds in this novel genuinely
do not settle, particularly on \texttt{exclusivity} and \texttt{adultery} in the
milieu relationships. The models almost never do. Pooling matched pairs, where the
annotator wrote \texttt{unclear} the models committed to a substantive label in
twelve cases out of thirteen; in the other direction they abstained on eleven of the
three hundred and seventeen fields where she had committed. This asymmetry shapes
what a corpus-scale run would produce: not a cautious annotation with visible gaps,
but a complete-looking table whose confident cells include the ones a human reader
would have left open. Reporting it alongside the field agreements also stops
abstention from becoming a cheap way to look accurate.

\subsection{What the measure cannot see}

Matching on participants is what makes the census score exact, and it also fixes what
the score can register.

\emph{Collective relationships are the residual problem.} Alias resolution settles
naming for named characters; it does not decide whether Laurette and her favourites
are one \texttt{hierarchy\_group} relationship or six dyads. That is a scholarly
decision, not noise, and under an exact key a legitimate disagreement about it costs
a false positive and a false negative at once. The gap between the dyadic and
collective strata in Table~\ref{tab:relation-results} is therefore partly a measure
of how each model resolved a question the codebook leaves open. Fixing a convention
--- a canonical identifier per milieu, with members listed in a separate field --- is
the first change we plan for the next version of the grid.

\emph{Direction is invisible.} The grid stores participants as a list, not as roles,
while several fields are directional: \texttt{reciprocity} = \texttt{one\_sided} does
not say who desires whom, and \texttt{adultery} = \texttt{adultery\_one} does not say
which party breaks the rule. A model that inverts the seducer and the seduced scores
a perfect match. For a project whose object is the asymmetry of desire this is a
substantive gap, and the reason a \texttt{direction} field is the second planned
change.

\emph{A strict metric is only as good as the label vocabulary.} Exact matching makes
the score reproducible, but it also means that a variant spelling of a label counts
as a disagreement. One misspelling in the gold file --- \texttt{asymetrical} for the
codebook's \texttt{asymmetrical} --- silently drove agreement on
\texttt{reciprocity} from 0.70 down to 0.25, turning a working field into an apparent
failure. The scoring pipeline now validates every label in both files against the
codebook's declared vocabulary before any comparison is made, and we report this
because it is the kind of error that survives peer review unnoticed.

\emph{The gold standard is one novel, and the counts are small.} Ten matched pairs
for the best model means a single relationship moves CQ by ten points, and the
ordering of Opus, Terra, and Sonnet in Table~\ref{tab:model-results} should not be
read as a separation. One of the fourteen gold relationships (Marion and the monk at
Solesmes) carries the annotator's own \texttt{exclude} flag and is nonetheless
counted against the models here; treating it as an ignore region, as detection
benchmarks do, would raise every recall figure. And we have no human ceiling yet:
fields where two human annotators would agree at 0.6 cannot be used to convict a
model at 0.6, so double annotation of a subset is the measurement that would make
Table~\ref{tab:field-results} fully interpretable.

\subsection{What we carry forward}

On this pilot, Claude Opus 5 is the model we take to the rest of the corpus: it is
the only one that annotates the novel at the granularity the codebook asks for,
milieux included, and it pays for that coverage with a precision the others never
have to risk. The ranking rests on one novel and fourteen relationships and should be
revisited as the gold standard grows.

More important than the ranking is the shape of the result. The census and the
descriptive half of the grid --- who is involved, in what gender configuration, in
what shape of group, in what kind of bond, with how much narrative room --- are
reliable enough to run over four hundred novels and to support distributional claims.
The normative half --- what the relation owes to whom, and what the novel makes of it
--- is not, and we do not intend to publish corpus-wide figures on
\texttt{framing} on the strength of an automatic pass. The pipeline's role in this
project is accordingly a division of labour rather than a replacement: it produces
the census and the descriptive coding at scale, and it directs human reading towards
the relationships where the interpretive fields are actually at stake.
